%

\subsection{Problem Cramera-Castillona} % lang-pl
\subsection{Problema di Cramer-Castillon} % lang-it

Poniższy problem będzie mieć ciekawą historię. % lang-pl
Il seguente problema avrà una storia interessante. % lang-it

\index{zadanie!Cramera-Castillona}
\index{problema!di Cramer-Castillon}
\begin{problem}
\label{problem_cramera_castillona}%
    Dane jest $n$ punktów oraz okrąg. % lang-pl
    Znaleźć taki wielokąt wpisany w~okrąg, że na każdym jego boku (lub jego przedłużeniu) znajduje się dokładnie jeden z~danych punktów. % lang-pl
    %
    Sono dati $n$ punti e un cerchio. % lang-it
    Trovare un poligono inscritto nel cerchio tale che su ciascun lato (o sul suo prolungamento) si trovi esattamente uno dei punti assegnati. % lang-it
\end{problem}
% TODO: https://ems.press/content/serial-article-files/45288 ładny obrazek, związek z Urquhartem

Szczególny przypadek $n = 3$ współliniowych punktów zainteresuje Pappusa. % lang-pl
Il caso particolare di $n = 3$ punti allineati interesserà Pappo. % lang-it
\index[persons]{Pappus}% % lang-pl
\index[persons]{Pappo}% % lang-it
Pewien nieznany, ale stary geometra przedstawi problem dla dowolnych trzech punktów Gabrielowi Cramerowi, który w~1742 przekaże go dalej do Giovanniego Salveminiego\footnote{Giovanni Francesco Mauro Melchiorre Salvemini di Castiglione przyjmie w pewnym roku nazwisko od swojego miejsca urodzenia, Castiglione del Valdarno w Toskanii.}. % lang-pl
Un geometra ignoto, ma antico, presenterà il problema per tre punti arbitrari a Gabriel Cramer, che nel 1742 lo trasmetterà a Giovanni Salvemini\footnote{Giovanni Francesco Mauro Melchiorre Salvemini di Castiglione assunse a un certo punto il cognome dal suo luogo di nascita, Castiglione del Valdarno in Toscana.}. % lang-it
\index[persons]{Castiglione, Giovanni}%
\index[persons]{Cramer, Gabriel}%
Ten znajdzie geometryczne rozwiązanie już po śmierci Cramera, w~1776 roku; wyczyn powtórzą Euler (1783 rok), Lagrange, Malfatti, Lhuilier, Servois, Poncelet % lang-pl
Egli troverà una soluzione geometrica dopo la morte di Cramer, nel 1776; il risultato sarà poi ottenuto nuovamente da Eulero (1783), Lagrange, Malfatti, L'Huillier, Servois e Poncelet % lang-it
\index[persons]{Lagrange, Joseph Louis}%
\index[persons]{Malfatti, Gian Francesco}%
\index[persons]{L'Huillier, Simon Antoine Jean}%
\index[persons]{Servois, François-Joseph}%
\index[persons]{Poncelet, Jean-Victor}%
% https://cms.math.ca/wp-content/uploads/crux-pdfs/Crux_v9n04_Apr.pdf Crux Mathematicorum, Vol 9, No 4, page 125
i szesnastoletni Annibale\footnote{Carnot uzna, że Ottajano (miejsce urodzenia Giordana) jest tytułem szlacheckim, a~nie nazwą miejscowości; w~swoich publikacjach nazwie młodego matematyka właśnie Ottajano. Niestety, ale później będzie tak określany także w~kolejnych pracach naukowych.} Giordano. % lang-pl
e dal sedicenne Annibale\footnote{Carnot ritenne che Ottajano (luogo di nascita di Giordano) fosse un titolo nobiliare anziché il nome di una località; nelle sue pubblicazioni chiamò quindi il giovane matematico Ottajano. Purtroppo questa denominazione verrà ripresa anche in lavori successivi.} Giordano. % lang-it
\index[persons]{Giordano, Annibale}%
Carnot uprości analityczne rozwiązanie Lagrange'a i~uogólni je do $n \ge 3$ punktów. % lang-pl
Carnot semplificherà la soluzione analitica di Lagrange e la generalizzerà al caso di $n \ge 3$ punti. % lang-it
\index[persons]{Carnot, Lazare}%

% TODO: https://en.wikipedia.org/wiki/Cramer-Castillon_problem

%